% 群论小抄模板
% 编译方式：XeLaTeX

\documentclass[landscape]{article}

% 中文支持
\usepackage{ctex}

% 页面设置：A3纸，页边距2mm
\usepackage[
    a3paper,
    landscape,
    top=2mm,
    bottom=2mm,
    left=2mm,
    right=2mm
]{geometry}
\usepackage{ytableau}
% 多栏排版
\usepackage{multicol}

% 删除页眉页脚
\pagestyle{empty}

% 设置字号为八号（约5pt）
\renewcommand{\normalsize}{\fontsize{5pt}{8pt}\selectfont}

% 设置行距为固定值8磅
\setlength{\baselineskip}{8pt}
\linespread{1}

% 分栏设置：6栏，间隔0.5字符（约2.5pt），加分隔线；关闭均衡分栏
\setlength{\columnsep}{2.5pt}  % 0.5字符间隔
\setlength{\columnseprule}{0.4pt}  % 分隔线宽度
\raggedcolumns  % 逐栏填满后再换栏，避免自动均衡

% 减少段落间距
\setlength{\parskip}{0pt}
\setlength{\parsep}{0pt}
\setlength{\parindent}{0pt}

% 紧凑的列表环境
\usepackage{enumitem}
\setlist{nosep, leftmargin=*, labelindent=0pt}

% 数学公式支持
\usepackage{amsmath, amssymb, amsfonts}

% 颜色支持（可选，用于标题高亮）
\usepackage{xcolor}

% 紧凑的section标题
\usepackage{titlesec}
\titleformat{\section}{\bfseries\normalsize}{\thesection.}{0.3em}{}
\titlespacing{\section}{0pt}{2pt}{1pt}
\titleformat{\subsection}{\bfseries\small}{\thesubsection.}{0.3em}{}
\titlespacing{\subsection}{0pt}{1pt}{0.5pt}

%======================================
% 文档开始
%======================================
\begin{document}
\normalsize

\begin{multicols*}{6}

% ===== 在此处添加你的内容 =====
\section*{群的定义与基本性质}
\textbf{群}：给定集合$G$及其上的二元运算（乘法），若满足以下四条，则称$(G,\cdot)$为群：封闭性、结合律、存在单位元、每个元素存在逆元。\\
\textbf{封闭性}：任取$f,g\in G$，其积$fg$仍属于$G$。\\
\textbf{结合律}：任取$f,g,h\in G$，先后相乘的结果与加括号方式无关。\\
\textbf{单位元}：存在且唯一的元素$e\in G$，使任取$f\in G$都有$ef=fe=f$。\\
\textbf{逆元}：任取$f\in G$，存在且唯一的元素$f^{-1}\in G$，使$f^{-1}f=ff^{-1}=e$。\\
\textbf{意义}：群论用于刻画与研究对称性。

\section*{例子}
\begin{itemize}
  \item $S_n$置换群：$n$个元素的置换群。
  \item $D_n$二面体群：正$n$边形的对称群。
  \item $A_n$交错群：是$S_n$的子群，且由偶置换构成。
\end{itemize}
$n$阶\textbf{置换群} $S_n$, 又称$n$阶\textbf{对称群}。
将$n$个元素的集合 $X=\{1, 2, \cdots, n\}$ 映为自身的置换为
\[
P = \begin{pmatrix} 
1 & 2 & \cdots & n \\ 
m_1 & m_2 & \cdots & m_n 
\end{pmatrix},
\]
其中 $m_1 m_2 \cdots m_n$ 是 $1, 2, \cdots, n$ 的任意排列, 
$P$ 表示把 1 映为 $m_1$, $n$ 映为 $m_n$ 的映射.。
显然置换只与每列的相对符号有关, 与第一行符号的顺序无关, 如
\[
\begin{pmatrix} 
1 & 2 & 3 & 4 \\ 
4 & 2 & 1 & 3 
\end{pmatrix} 
= 
\begin{pmatrix} 
4 & 2 & 3 & 1 \\ 
3 & 2 & 1 & 4 
\end{pmatrix}.
\]

\noindent 定义两个置换 $P'$ 和 $P$ 的乘积 $P'P$, 为先实行置换 $P$, 再实行置换 $P'$, 如
\[
\begin{pmatrix} 
1 & 2 & 3 \\ 
2 & 1 & 3 
\end{pmatrix}
\begin{pmatrix} 
1 & 2 & 3 \\ 
3 & 2 & 1 
\end{pmatrix}
=
\begin{pmatrix} 
1 & 2 & 3 \\ 
3 & 1 & 2 
\end{pmatrix}.
\]
容易看出在这乘法定义下, 全部 $n$ 阶置换构成 $S_n$ 群. $S_n$ 群共有 $n!$ 个元素

$D_3$ 群元素定义：$e$(单位元), $a$(关于轴 1 反射), 
$b$(关于轴 2 反射), $c$(关于轴 3 反射), $d$(绕垂直中心轴转 $120^\circ$), $f$(绕垂直中心轴转 $240^\circ$)。
\begin{tabular}{|c|c c c c c c|}
    \hline
      & $\boldsymbol{e}$ & $\boldsymbol{d}$ & $\boldsymbol{f}$ & $\boldsymbol{a}$ & $\boldsymbol{b}$ & $\boldsymbol{c}$ \\
    \hline
    $\boldsymbol{e}$ & $e$ & $d$ & $f$ & $a$ & $b$ & $c$ \\
    $\boldsymbol{d}$ & $d$ & $f$ & $e$ & $c$ & $a$ & $b$ \\
    $\boldsymbol{f}$ & $f$ & $e$ & $d$ & $b$ & $c$ & $a$ \\
    $\boldsymbol{a}$ & $a$ & $b$ & $c$ & $e$ & $d$ & $f$ \\
    $\boldsymbol{b}$ & $b$ & $c$ & $a$ & $f$ & $e$ & $d$ \\
    $\boldsymbol{c}$ & $c$ & $a$ & $b$ & $d$ & $f$ & $e$ \\
    \hline
\end{tabular}
\par

\section*{群的元素/阶/可交换性}
\begin{itemize}
  \item 群元素可以是数，也可以是变换（矩阵）、置换等。
  \item \textbf{有限群/无限群}：$|G|<\infty$ 为有限群，否则为无限群；\textbf{群的阶}为$|G|$（也记$n(G)$）。
  \item 有限群的乘法规则常用\textbf{乘法表}描述；无限群常用\textbf{公式}描述。
  \item \textbf{阿贝尔群（可交换群）}：对任意$f,g\in G$，有$fg=gf$。
\end{itemize}

\section*{群元素的参数标记}
群元素可用指标/参数$\alpha$标记为$g_{\alpha}$（$\alpha$可取离散或连续范围；也可由多个参数共同标记）。

\section*{重排定理（Rearrangement）}
\begin{itemize}
  \item 对任意固定$u\in G$，左乘集合$uG=\{ug\mid g\in G\}$与$G$相同；右乘集合$Gu=\{gu\mid g\in G\}$与$G$相同。
  \item 取逆集合$G^{-1}=\{g^{-1}\mid g\in G\}$与$G$相同。
  \item 等价表述：在乘法表中，每一行/列都是$G$元素的一个排列；每个元素在每一行/列中恰出现一次。
\end{itemize}
\noindent\textcolor{red}{\rule{\linewidth}{1pt}}
\section*{子群}
\begin{itemize}
  \item \textbf{定义}：$H\subseteq G$且在$G$的同一运算下构成群，则称$H$为$G$的子群（记$H \subset  G$）。
  \item \textbf{判别（非空子集）}：$H$为子群$\Leftrightarrow$（1）$h_1,h_2\in H\Rightarrow h_1h_2\in H$；（2）$h\in H\Rightarrow h^{-1}\in H$。
  \item \textbf{平凡子群}：$\{e\}$与$G$；其余称固有子群。
\end{itemize}

\section*{循环群与元素的阶}
\begin{itemize}
  \item 任意元素$a\in G$可生成循环子群$\langle a\rangle=\{a^k\mid k\in\mathbb{Z}\}$。
  \item 若存在最小正整数$k$使$a^k=e$，则称$a$的\textbf{阶}为$k$，且$\langle a\rangle$为$k$阶循环群（有限情形）。
  \item 循环群为阿贝尔群。
\end{itemize}

\section*{陪集}
\begin{itemize}
  \item 若{\color{red}子群}$H\subset G$，$g\in G$，\textbf{左陪集}$gH=\{gh\mid h\in H\}$；\textbf{右陪集}$Hg=\{hg\mid h\in H\}$。
  \item 任一陪集的元素数目与子群$H$相等。
\end{itemize}

\section*{陪集定理与拉格朗日定理}
\begin{itemize}
  \item \textbf{陪集定理}：任意两个左陪集要么相同，要么不相交；右陪集同理。
  \item \textbf{拉格朗日定理}：若$G$为有限群、$H\subset G$，则$|H|\mid |G|$，且$|G|=|H|\,[G\!:\!H]$（$[G\!:\!H]$为左陪集个数）。
  \item \textbf{推论}：若$|G|$为素数，则$G$无非平凡子群（只有$\{e\}$与$G$）。
  \item \textbf{分割}：有限群$G$可按某子群$H$的左（或右）陪集分割为互不相交子集。
\end{itemize}

$D_3$ 有子群 $H_1=\{e, a\}$, $H_2=\{e, b\}$, $H_3=\{e, c\}$ 和 $H_4=\{e, d, f\}$。
\noindent $D_3$ 可按 $H_1$ 分成左陪集串， $H_1=\{e, a\}, bH_1=\{b, f\}, cH_1=\{c, d\}$。\\
也可按 $H_4$ 分成右陪集串， $H_4=\{e, d, f\}, H_4a=\{a, b, c\}$。
\noindent\textcolor{red}{\rule{\linewidth}{1pt}}
\section*{共轭与共轭类}
\begin{itemize}
  \item \textbf{共轭}：$f,h\in G$，若$\exists g\in G$使$gfg^{-1}=h$，则称$h$与$f$共轭，记$h\sim f$。
  \item \textbf{性质}：自反（$f\sim f$）、对称、传递；故“共轭”是等价关系。
  \item \textbf{共轭类}：与$ f $共轭的全体元素构成$ f $的共轭类$[f]=\{gfg^{-1}\mid g\in G\}$。
  \item \textbf{基本事实}：单位元自成一类；阿贝尔群每个元素自成一类；同一共轭类中元素的阶相同。
  \item \textbf{分割}：不同共轭类互不相交；群可按共轭类分割（各类大小不一定相同）。
\end{itemize}

\section*{中心化子（对易子群）}
\begin{itemize}
  \item \textbf{中心化子（对易子群）}：$C_G(g)=\{h\in G\mid hg=gh\}=\{h\in G\mid hgh^{-1}=g\}$，且$C_G(g)\subset G$。
  \item \textbf{共轭类大小}：若$G$有限，则$|[g]|=[G:C_G(g)]=\dfrac{|G|}{|C_G(g)|}$。
  \item \textbf{定理1.4}：有限群中任一共轭类的元素个数是群阶的因子（即$|[g]|\mid |G|$）。
\end{itemize}

\section*{共轭子群}
\begin{itemize}
  \item \textbf{定义}：$H,K\subset G$，若$\exists g\in G$使$K=gHg^{-1}=\{ghg^{-1}\mid h\in H\}$，则称$H$与$K$共轭。
  \item \textbf{性质}：对称、传递（因此子群可按共轭关系分成共轭子群类）。
\end{itemize}

\section*{不变子群（正规子群）}
\begin{itemize}
  \item \textbf{定义}：$H\subset G$，若$\forall g\in G$有$gHg^{-1}=H$，则称$H$为$G$的不变子群（正规子群），记$H\trianglelefteq G$。
  \item \textbf{等价表述}：$\forall g\in G,\forall h\in H$，有$ghg^{-1}\in H$（即$H$包含其任一元素的整个共轭类）。
  \item \textbf{阿贝尔群}：若$G$为阿贝尔群，则$G$的任意子群都是不变子群。
  \item \textbf{定理1.5（共轭重排）}：若$H\trianglelefteq G$且$f\in G$固定，则映射$h\mapsto fhf^{-1}$是$H\to H$的双射（在$h$取遍$H$时，$fhf^{-1}$恰给出$H$全部元素一次）。
\end{itemize}

$D_3$ 群的非平庸子群
\begin{itemize}
    \item \textbf{2阶循环群：} $H_1=\{e, a\},\; H_2=\{e, b\},\; H_3=\{e, c\}$
    \item \textbf{3阶循环群：} $H_4=\{e, d, f\}$
\end{itemize}
$\{H_1, H_2, H_3\}$是共轭子群类；
$H_4$ 是自共轭子群（不变子群/正规子群），即 $g^{-1}H_4 g = H_4$。

\section*{不变子群的陪集与商群}
\begin{itemize}
  \item 若$H\trianglelefteq G$，则$\forall g\in G$有\textbf{左右陪集重合}：$gH=Hg$（不再区分左/右陪集）。
  \item \textbf{陪集乘法}：$(g_iH)(g_jH)=g_ig_jH$（良定义依赖$H\trianglelefteq G$）。
  \item \textbf{商群}：以陪集集合$G/H=\{gH\mid g\in G\}$为元素、以上述陪集乘法为运算，则$G/H$构成群，称为$H$的商群。
\end{itemize}
例子：商群D3/H4= Z2

\noindent\textcolor{red}{\rule{\linewidth}{1pt}}
\section*{同构（Isomorphism）}
\begin{itemize}
  \item \textbf{定义1.10}：若存在双射$\Phi:G\to F$满足$\Phi(g_ig_j)=\Phi(g_i)\Phi(g_j)$，则称$G$与$F$同构，记$G\cong F$；$\Phi$称同构映射。
  \item \textbf{保结构}：同构映射将单位元映到单位元，将逆元映到逆元：$\Phi(e_G)=e_F$，$\Phi(g^{-1})=\Phi(g)^{-1}$。
  \item \textbf{逆映射}：若$\Phi$为同构，则$\Phi^{-1}$存在且也是同构；因此$G\cong F\Rightarrow F\cong G$。
  \item \textbf{观点}：同构的两个群可视为“同一群的不同表示”（群结构相同）。
  \item \textbf{例子}：$D_3\cong S_3$，空间反演群和$Z_2$，互为共轭的两个子群。
\end{itemize}

\section*{同态（Homomorphism）}
\begin{itemize}
  \item \textbf{定义1.11}：映射$\Phi:G\to F$若保持乘法：$\Phi(g_ig_j)=\Phi(g_i)\Phi(g_j)$，则称$G$到$F$的同态（若$\Phi$为满射，也常称满同态）。
  \item \textbf{可放宽}：亦可不要求$\Phi$满射；若不特别说明，可按需要区分“同态/满同态”。
  \item \textbf{关系}：同构是特殊同态（同态$\Phi$若为双射$\Leftrightarrow$同构）。
  \item \textbf{基本性质}：$\Phi(e_G)=e_F$，$\Phi(g^{-1})=\Phi(g)^{-1}$。
  \item \textbf{一般性}：同态可多对一；$G$同态到$F$不必推出$F$同态到$G$。
  \item \textbf{显然同态}：任意群$G$到平凡群$Z_1=\{e\}$存在同态（通常略去）。
\end{itemize}

\section*{同态核}
\begin{itemize}
  \item \textbf{同态核}：$\ker\Phi=\{g\in G\mid \Phi(g)=e_F\}$。
  \item \textbf{性质}：$\ker\Phi\subset G$且$\ker\Phi\trianglelefteq G$（同态核是不变子群）。
  \item \textbf{同态核定理}：(1).$\ker\Phi\trianglelefteq G$；(2)$G/\ker\Phi\cong F$。
  \item \textbf{例子}：$D_3$ 到 $Z_2$ 的同态，$\ker\Phi=\{e, d, f\}$。
\end{itemize}

\section*{自同构与自同构群}
\begin{itemize}
  \item \textbf{自同构映射}：群$G$到自身的映射$\nu:G\to G$称为$G$的自同构映射。
  \item \textbf{自同构群}：所有自同构映射在复合运算下构成群，记$\operatorname{Aut}(G)$。
\end{itemize}

\section*{内自同构映射与内自同构群}
\begin{itemize}
  \item \textbf{内自同构}：由$u\in G$诱导的共轭映射$\mu_u(g)=ugu^{-1}$称为内自同构映射。
  \item \textbf{内自同构群}：全体内自同构构成群$\operatorname{Inn}(G)$，且$\operatorname{Inn}(G)\trianglelefteq \operatorname{Aut}(G)$。
\end{itemize}

\section*{群的中心}
\begin{itemize}
  \item \textbf{中心}：$Z(G)=\{z\in G\mid \forall g\in G,\ zg=gz\}$；若$G$为阿贝尔群，则$Z(G)=G$。
  \item \textbf{性质}：$Z(G)\trianglelefteq G$。
  \item \textbf{结论}：$G/Z(G)\cong \operatorname{Inn}(G)$。
\end{itemize}

\section*{外自同构群}
\begin{itemize}
  \item \textbf{外自同构}：不是内自同构的自同构称为外自同构（outer automorphism）。
  \item \textbf{外自同构群}：$\operatorname{Out}(G)\equiv \operatorname{Aut}(G)/\operatorname{Inn}(G)$（因$\operatorname{Inn}(G)\trianglelefteq \operatorname{Aut}(G)$）。
\end{itemize}
\noindent\textcolor{red}{\rule{\linewidth}{1pt}}
\section*{变换群（置换群）}
\begin{itemize}
  \item \textbf{变换/置换}：非空集合$X$到自身的一一对应（双射）$f:X\to X$。
  \item \textbf{置换乘法（定义1.15）}：对$X$上的置换$f,g$，定义$fg$为先做$g$后做$f$，即$fg(x)=f(g(x))$。
  \item \textbf{完全对称群}：$X$上全体置换在上述乘法下成群，记为$S_X$；单位元为恒等置换$e$，逆元为逆置换$f^{-1}$。
  \item 若$|X|=n$，则$S_X$记为$S_n$，且$|S_n|=n!$。
  \item $S_X$的任意子群称为$X$上的\textbf{变换群/置换群/对称群}。
\end{itemize}

\section*{凯莱定理（Cayley）}
\begin{itemize}
  \item \textbf{定理1.7}：任意群$G$同构于$S_G$（以$G$为底集合的完全对称群）的一个子群；特别地，若$|G|=n$，则$G$同构于$S_n$的一个子群。
\end{itemize}

\section*{群作用下的等价、轨道与不变子集}
\begin{itemize}
  \item 设$G$为$X$的变换群。定义$X$上关系：$x\sim y \iff \exists g\in G,\ gx=y$；则$\sim$具对称性与传递性（通常也视为等价关系）。
  \item \textbf{轨道}：含$x$的$G$轨道为$\mathcal{O}_G(x)=\{gx\mid g\in G\}$（即与$x$等价的全体点）。
  \item \textbf{迷向子群/稳定子群}：$G_x=\{h\in G\mid hx=x\}$，是$G$的子群（使$x$不动的全部变换）。
  \item \textbf{轨道--陪集对应}：含$x$的轨道上的点与$G_x$的左陪集一一对应；
  \item \textbf{推论（有限情形）}：若$|G|=n$且$|G_x|=n(G_x)$，则$|\mathcal{O}_G(x)|=[G:G_x]=\dfrac{n}{n(G_x)}$。
  \item \textbf{$G$不变子集}：$Y\subseteq X$称$G$不变，若$\forall y\in Y,\forall g\in G,\ g(y)\in Y$；每个轨道都是$G$不变的，轨道的并也是$G$不变的。
  \item \textbf{子集的稳定子群}：对任意$Y\subseteq X$，可取$H=\{g\in G\mid g(Y)=Y\}$，则$H\subset G$且$Y$对$H$不变（至少平凡子群$\{e\}$总可使$Y$不变）。
\end{itemize}
\noindent\textcolor{red}{\rule{\linewidth}{1pt}}
\section*{直积群（Direct Product）}
\begin{itemize}
  \item 设$G_1,G_2$为群。直积群$G=G_1\otimes G_2$的元素可写为$g_{\alpha\beta}=g_{1\alpha}g_{2\beta}$，其中$g_{1\alpha}\in G_1,\,g_{2\beta}\in G_2$，且视作可交换：$g_{1\alpha}g_{2\beta}=g_{2\beta}g_{1\alpha}$。
  \item \textbf{乘法}：$g_{\alpha\beta}g_{\alpha'\beta'}=(g_{1\alpha}g_{1\alpha'})(g_{2\beta}g_{2\beta'})$（即两分量分别相乘再组合）。
  \item \textbf{单位元/逆元}：$e=e_1e_2$；$(g_{1}g_{2})^{-1}=g_1^{-1}g_2^{-1}$。
  \item 任何一个群乘以其它群的单位元，同构于原来的那个群。
  设$e_1,e_2$分别为$G_1,G_2$的单位元，则
  $
    F_1 = \{ g_{1\alpha} e_2 \mid g_{1\alpha} \in G_1 \} \cong G_1,\quad
    F_2 = \{ e_1 g_{2\beta} \mid g_{2\beta} \in G_2 \} \cong G_2
  $。
  显然，$F_1$和$F_2$都是$G = G_1 \otimes G_2$的子群，可以写为$G = F_1 \otimes F_2$。
  \item \textbf{判别}：若群$G$有子群$G_1,G_2$满足（i）任意$g\in G$可唯一表示为$g=g_1g_2$（$g_i\in G_i$）；（ii）$g_1g_2=g_2g_1$（$\forall g_1\in G_1,g_2\in G_2$），则$G$是$G_1$与$G_2$的直积：$G=G_1\otimes G_2$。
  \item \textbf{性质}：若$G=G_1\otimes G_2$，则$G_1\cap G_2=\{e\}$，且$G_1\trianglelefteq G,\ G_2\trianglelefteq G$。
  \item \textbf{商群}：$(G_1\otimes G_2)/G_1\cong G_2,\quad (G_1\otimes G_2)/G_2\cong G_1$。
\end{itemize}

\section*{半直积群（Semidirect Product）}
\begin{itemize}
  \item 给定群$G_1,G_2$及同态$\Phi:G_2\to \operatorname{Aut}(G_1)$（记$\Phi(g_2)=\varphi_{g_2}$）。
  \item \textbf{元素表示}：$G=G_1\rtimes_\Phi G_2$的元素可唯一写为$(g_1,g_2)$（或有序乘积$g_1g_2$）。
  \item \textbf{乘法}：$(g_1,g_2)(g_1',g_2')=(g_1\,\varphi_{g_2}(g_1'),\ g_2g_2')$。
  \item \textbf{单位元}：$(e_1,e_2)$。
  \item \textbf{逆元}：$(g_1,g_2)^{-1}=\big(\varphi_{g_2^{-1}}(g_1^{-1}),\ g_2^{-1}\big)$。
  \item \textbf{结构特点}：一般不要求$g_1g_2=g_2g_1$；通常有$G_1\trianglelefteq G$，而$G_2$未必正规。
  \item \textbf{顺序不可交换}：一般$G_1\rtimes G_2 \not\cong G_2\rtimes G_1$（构造依赖作用$\Phi$）。
  \item \textbf{退化为直积}：若$\Phi$为平凡作用（$\forall g_2,\ \varphi_{g_2}=\mathrm{id}$），则$G_1\rtimes_\Phi G_2 \cong G_1\otimes G_2$（此时两分量可交换）。
  \item \textbf{判别}：$G$为$G_1$与$G_2$的半直积$\Rightarrow$ $G_1\trianglelefteq G$，$G_1\cap G_2=\{e\}$，且每个$g\in G$可唯一写为$g=g_1g_2$（$g_i\in G_i$）。
  \item \textbf{重要观点}：若$H\trianglelefteq G$，则常把$G$看成$H$与$G/H$的半直积：$G\cong H\rtimes (G/H)$。
\end{itemize}
\noindent\textcolor{red}{\rule{\linewidth}{1pt}}
\section*{线性空间（向量空间）}
\begin{itemize}
  \item \textbf{定义2.1}：数域$K$上的向量空间$V$在加法与数乘下满足：加法交换/结合、零元$0$、加法逆元$-x$；数乘单位$1\cdot x=x$、结合律$(ab)x=a(bx)$、分配律$a(x+y)=ax+ay$、$(a+b)x=ax+bx$。
  \item 视加法为“群乘法”，则$(V,+)$是阿贝尔群。
\end{itemize}

\section*{线性变换}
\begin{itemize}
  \item \textbf{定义2.2}：线性变换$A:V\to V$满足$A(ax+y)=aA(x)+A(y)$（$a\in K$）。
  \item 线性变换可定义数乘、加法与复合乘法：$(aA)(x)=aA(x)$，$(A+B)(x)=A(x)+B(x)$，$(AB)(x)=A(B(x))$。
  \item 若$A$为双射，则存在逆线性变换$A^{-1}$。
  \item \textbf{有限维}：$V$至多$n$个线性无关向量$\Rightarrow$ $V$为$n$维；可取基$(e_1,\dots,e_n)$，任意$x=\sum_i x_i e_i$（坐标列向量）。
  \item \textbf{矩阵表示}：线性变换在基下由$n\times n$矩阵表示；行列式$\neq 0$时为非奇异（可逆）变换。
\end{itemize}

\section*{一般线性群}
\begin{itemize}
  \item \textbf{$GL(n,\mathbb{C})$}：{\color{red}全体}$n\times n$复矩阵中{\color{red}行列式非零者}在矩阵乘法下构成群。
  \item \textbf{复数域的一般线性群}：$n$维复向量空间$V$上全体非奇异线性变换在复合下构成群，称为$GL(V,\mathbb{C})$（与$GL(n,\mathbb{C})$同构）；单位元为恒等变换（单位矩阵），逆元为逆变换（逆矩阵）。
\end{itemize}

\section*{群表示}
\begin{itemize}
  \item \textbf{线性表示}：群$G$到线性变换群$L(V,\mathbb{C})$的{\color{red}同态}$A:G\to L(V,\mathbb{C})$称为$G$在$V$上的（复）线性表示；$\dim V=n$称为$n$维表示。等价地（选定基后）$A:G\to GL(n,\mathbb{C})$为同态。
  \item \textbf{保单位元/逆元}：$A(e)=I$，$A(g^{-1})=A(g)^{-1}$。
  \item \textbf{为何要求非奇异}：若某$g$映为奇异矩阵（$\det=0$），由同态与重排定理将迫使所有$A(h)$奇异，与$A(e)=I$矛盾；故表示矩阵应可逆。
  \item \textbf{忠实表示}：若$A$为{\color{red}同构}（等价于$\ker A=\{e\}$），则称$A$为忠实表示。
  \item \textbf{平凡表示}：$1$维表示$A(g)=1$（或$A(g)=I_{1\times1}$）对任意$g\in G$。
\end{itemize}

\section*{由群作用诱导的表示（函数空间）}
\begin{itemize}
  \item 若$G$作用在空间$X$上，则可在函数空间上定义表示（拉回）：$(A(g)\psi)(x)=\psi(g^{-1}x)$。
  \item 该定义满足同态性：$A(g_1g_2)=A(g_1)A(g_2)$（从而给出$G$的一种表示构造方法）。
\end{itemize}

\section*{一维表示与可约性（概念）}
\begin{itemize}
  \item \textbf{一维表示}：本质上是群同态$A:G\to \mathbb{C}^{\times}$（把群元映到可逆标量）。
  \item \textbf{可约表示（直和分解）}：若表示空间可分解为非平凡$G$不变子空间的直和，则称表示可约；否则不可约。
\end{itemize}

\section*{由已知表示构造高维表示（多项式/张量观点）}
\begin{itemize}
  \item 已知$G$在$V$上的表示$A$后，可在$V$的多项式函数空间（如二次齐次多项式所张成的空间）上用拉回$(A(g)\phi)(v)=\phi(A(g)^{-1}v)$得到更高维表示。
  \item 这类构造常产生可约表示，进而可分解为直和（如写作$V\simeq V_1\oplus V_2$，对应表示$A\simeq A_1\oplus A_2$）。
  \item 若群G保持系统的
  哈密顿量不变，则我们
  哈密顿量不变，则我们
  可用某个能量本征值对
  应的简并子空间的所有
  本征波函数为基底 构
  造群G的矩阵表示，矩
  阵的维数等于该能级的
  简并度。
\end{itemize}

\section*{从不变子空间读出表示矩阵}
\begin{itemize}
  \item 若$W$是表示空间中一个$G$不变子空间，取基$(\psi_1,\dots,\psi_m)$，则对任意$g\in G$有展开
  \[
  A(g)\psi_i=\sum_{j=1}^m \psi_j\,A_{ji}(g),
  \]
  从而得到矩阵表示$A(g)=(A_{ji}(g))\in GL(m,\mathbb{C})$。
  \item 同态性对应矩阵乘法：$A(g_1g_2)=A(g_1)A(g_2)$。
\end{itemize}
\noindent\textcolor{red}{\rule{\linewidth}{1pt}}
\section*{等价表示（相似变换）}
\begin{itemize}
  \item \textbf{定义2.6}：若$A$是$G$在$V$上的表示，$X$为可逆线性变换，则
  \[
  A'(g)=XA(g)X^{-1}
  \]
  也是表示，称$A'$与$A$等价（equivalent）。
  \item 等价表示维数相同；维数相同的表示未必等价。
  \item 相似变换本质是更换表示空间的基。
\end{itemize}

\section*{可约/完全可约/不可约表示}
\begin{itemize}
  \item \textbf{（可约）}：表示$A$在$V$上若存在$G$不变的真子空间$W$（$0\neq W\neq V$），则$A$可约；否则不可约。
  \item \textbf{矩阵判别（可约）}：$A$可约$\Leftrightarrow$存在等价表示使全部$A(g)$在同一基下呈共同的块上（三）角形式
  \[
  A(g)=\begin{pmatrix} C(g)&N(g)\\ 0&B(g)\end{pmatrix}.
  \]
  \item \textbf{直和}：$V=W\oplus W'$指$V=W+W'$且$W\cap W'=\{0\}$；任意$x\in V$可唯一写为$x=y+z$（$y\in W,z\in W'$）。
  \item \textbf{（完全可约）}：若$V=W\oplus W'$且$W,W'$均为$G$不变子空间，则$A$完全可约（是可约的特殊情形）。
  \item \textbf{矩阵判别（完全可约）}：存在等价表示使全部$A(g)$共同块对角化
  \[
  A(g)=C(g)\oplus B(g)=\begin{pmatrix} C(g)&0\\ 0&B(g)\end{pmatrix},
  \]
  即表示可写为直和（可包含重数）。{\color{red}一个完全可约表示可写为多个不可约表示的直和}，如果其中某些
  不可约表示有重复，还可以用其倍数来表示。
  \item \textbf{不可约}：表示空间不存在$G$不变的真子空间。
  \item \textbf{直接推论}：若$A$不可约，则不存在等价表示使\emph{所有} $A(g)$共同化为块上三角或块对角形式（否则将给出不变真子空间）。
\end{itemize}

\section*{内积空间与正交}
\begin{itemize}
  \item \textbf{内积}：在数域$K$上的线性空间$V$中，映射$(x|y)\in K$满足共轭对称$( x|y)=( y|x)^{*}$、对第二个变量线性、正定$( x|x)>0$（$x\neq0$），则称为内积。
  \item \textbf{长度/正交}：$|x|=\sqrt{( x|x)}$；若$( x|y)=0$则$x\perp y$。
  \item \textbf{正交归一基}：可取基$\{e_i\}$使$( e_i|e_j)=\delta_{ij}$。
\end{itemize}

\section*{幺正变换与幺正表示}
\begin{itemize}
  \item \textbf{幺正变换}：内积空间$V$上线性变换$U$若保持内积$(Ux|Uy)=(x|y)$，则$U$为幺正变换。
  \item \textbf{性质}：$U^{\dagger}U=UU^{\dagger}=I$，故$U^{\dagger}=U^{-1}$（矩阵形式：$U^{*t}U=I$）。
  \item \textbf{幺正/酉表示}：若表示$A:G\to GL(V)$的每个$A(g)$都是幺正变换，则称$A$为$G$的幺正表示；并有$A(g)^{\dagger}=A(g)^{-1}=A(g^{-1})$。
\end{itemize}

\section*{幺正表示的完全可约性}
\begin{itemize}
  \item \textbf{公式}：在正交归一基下，幺正表示矩阵满足$A(g)^{*t}=A(g)^{-1}=A(g^{-1})$（分量式：$A(g)_{ji}^{*}=[A(g)^{-1}]_{ij}=A(g^{-1})_{ij}$）。
  \item \textbf{完全可约性}：若幺正表示可约，则必完全可约（可用不变子空间$W$与其正交补$W^{\perp}$分解：$V=W\oplus W^{\perp}$）。
  \item \textbf{推论}：有限维幺正表示可分解为不可约幺正表示的直和（可带重数）：
  
  $A=\bigoplus_{p} m_p A^{p}$。
\end{itemize}
\noindent\textcolor{red}{\rule{\linewidth}{1pt}}
\section*{线性代数与结合代数}
\begin{itemize}
  \item \textbf{定义}：数域$K$上线性空间$R$若定义了乘法，并对$x,y,z\in R,a\in K$满足：封闭$xy\in R$；分配律$x(y+z)=xy+xz,(x+y)z=xz+yz$；数乘可提出$a(xy)=(ax)y=x(ay)$，则称$R$为（线性）代数。
  \item 若还满足结合律$(xy)z=x(yz)$，则称为\textbf{结合代数}。
\end{itemize}

\section*{群空间与群代数}
\begin{itemize}
  \item \textbf{群空间}：取复数域$\mathbb{C}$与群$G=\{g_\alpha\}$，令
  \[
  V_G=\left\{x=\sum_{\alpha} x_{\alpha}g_{\alpha}\mid x_\alpha\in\mathbb{C}\right\},
  \]
  并定义$(x+y)=\sum_\alpha(x_\alpha+y_\alpha)g_\alpha,\ a x=\sum_\alpha(ax_\alpha)g_\alpha$，则$V_G$为线性空间，$\{g_\alpha\}$为自然基。
  \item \textbf{群代数}：若$g_\alpha g_\beta=g_\gamma$，对$x=\sum_\alpha x_\alpha g_\alpha,\ y=\sum_\beta y_\beta g_\beta$定义乘法
  \[
  xy=\sum_{\alpha,\beta} x_\alpha y_\beta (g_\alpha g_\beta)\in V_G,
  \]
  该乘法与加法/数乘满足分配与结合，从而$V_G$成为$G$的结合代数（群代数）。
\end{itemize}

\section*{正则表示（Regular Representation）}
\begin{itemize}
  \item 以群代数$V_G$为表示空间。
  \item \textbf{左正则表示}：对$g\in G$定义线性变换$L(g)$满足$L(g)\,g_\alpha=gg_\alpha$（再线性延拓到$V_G$）；则$L(g_1g_2)=L(g_1)L(g_2)$，给出$G$的表示；若$|G|=n$则为$n$维表示。
  \item \textbf{忠实性}：由重排定理，$g_1\neq g_2\Rightarrow L(g_1)\neq L(g_2)$，故左正则表示忠实。
  \item \textbf{右正则表示}：定义$R(g)\,g_\alpha=g_\alpha g^{-1}$（线性延拓），同样满足$R(g_1g_2)=R(g_1)R(g_2)$，亦为$G$的表示。
  \item \textbf{矩阵形态（有限群）}：在自然基$\{g_\alpha\}$下，$L(g)$（或$R(g)$）对应的矩阵是置换矩阵：元素只取$0/1$，且每行每列恰有一个$1$（其余为$0$）。
  \item \textbf{约化（概念）}：正则表示一般是可约的；可通过相似变换将所有$L(g)$共同化为块对角形式，即分解为若干表示的直和。
\end{itemize}
\noindent\textcolor{red}{\rule{\linewidth}{1pt}}
\section*{有限群表示的幺正化}
\begin{itemize}
  \item \textbf{定理（有限群表示等价于幺正表示）}：若$G$为有限群，$A$为$G$在有限维内积空间$V$上的表示，则存在等价表示$A'(g)=X A(g) X^{-1}$使$A'$为幺正表示。
  \item \textbf{方法（平均内积）}：在原内积$(\cdot|\cdot)$下定义新内积
  \[
  \langle x|y\rangle=\frac{1}{|G|}\sum_{g\in G} (A(g)x\,|\,A(g)y),
  \]
  则$A(g)$在$\langle\cdot|\cdot\rangle$下成为幺正变换。
  \item \textbf{推论}：{\color{red}有限群在内积空间的表示若可约则必完全可约}。
\end{itemize}
\noindent\textcolor{red}{\rule{\linewidth}{1pt}}
\section*{Schur 引理（一）}
\begin{itemize}
  \item 设$A,B$分别为群$G$的$n_A$维、$n_B$维不可约表示。若存在矩阵$M$满足
  \[
  A(g)M=MB(g)\quad (\forall g\in G),
  \]
  则：（i）若$A,B$不等价，则$M=0$；（ii）若$M\neq0$，则$A,B$等价。
  \item \textbf{推论}：若$M\neq0$，则$M$为同构（可逆），且$n_A=n_B$。
\end{itemize}

\section*{Schur 引理（二）}
\begin{itemize}
  \item 设$A$为群$G$的有限维\textbf{不可约复表示}。若矩阵$M$与所有$A(g)$对易：
  \[
  A(g)M=MA(g)\quad (\forall g\in G),
  \]
  则$M=\lambda I$（$\lambda\in\mathbb{C}$）。
  \item \textbf{备注}：该结论依赖“复表示/代数闭域”；对\textbf{实}不可约表示不一定成立。
\end{itemize}
\noindent\textcolor{red}{\rule{\linewidth}{1pt}}
\section*{群函数}
\begin{itemize}
  \item \textbf{对应}：群代数（群空间）$V_G$中向量$x=\sum_i x_i g_i$可视为群函数$x(g_i)\equiv x_i$；反之给定函数值$x(g_i)$即可确定$x=\sum_i x(g_i)g_i$。
  \item \textbf{数乘/加法}：$(ax)(g_i)=a\,x(g_i)$；$(x+y)(g_i)=x(g_i)+y(g_i)$。
  \item \textbf{乘法（卷积型）}：
  \[
  (xy)(g_i)=\sum_j x(g_j)\,y(g_j^{-1}g_i).
  \]
  \item \textbf{一组基函数}：可取$n=|G|$个函数$g_i(\,g_j\,)=\delta_{ij}$作为群函数空间的基（对应群代数的自然基$g_i$）。
\end{itemize}

\section*{群函数的内积与正则表示}
\begin{itemize}
  \item 设$G=\{g_i\}$为$n$阶有限群。定义群函数内积
  \[
  \langle x|y\rangle=\frac{1}{n}\sum_{i=1}^{n} x(g_i)^{*}\,y(g_i).
  \]
  \item 特别地，对$\delta$基函数有$\langle g_i|g_j\rangle=\frac{1}{n}\delta_{ij}$，是正交完备的。
  \item 在该内积下，左正则表示$L(g)$为幺正：$\langle L(g)x|L(g)y\rangle=\langle x|y\rangle$。
  \item 若$A^{(r)}$为$n_r$维表示，则其矩阵元函数$A^{(r)}_{\mu\nu}(g_i)$（$\mu,\nu=1,\dots,n_r$）都是群函数；共$n_r^2$个。
\end{itemize}
\noindent\textcolor{red}{\rule{\linewidth}{1pt}}
\section*{正交性定理（矩阵元）}
\begin{itemize}
  \item 设$\{A^{(p)}\}$为有限群$G$的一组不等价不可约（幺正）表示，维数为$n_p$，则
  \[
  \sum_{i=1}^{n}\left[A^{(p)}_{\mu\nu}(g_i)\right]^{*}A^{(r)}_{\mu'\nu'}(g_i)
  =\frac{n}{n_p}\,\delta_{pr}\delta_{\mu\mu'}\delta_{\nu\nu'}.
  \]
  \item 等价写法（用上面的内积）：$\langle A^{(p)}_{\mu\nu}\,|\,A^{(r)}_{\mu'\nu'}\rangle=\frac{1}{n_p}\delta_{pr}\delta_{\mu\mu'}\delta_{\nu\nu'}$。
\end{itemize}

\section*{完备性与正则表示约化}
\begin{itemize}
  \item \textbf{完备性定理}：有限群$G$的全部不等价不可约（幺正）表示的矩阵元群函数$A^{(r)}_{\mu\nu}(g_i)$在群函数空间中构成完备系；任意群函数可关于它们展开。
  \item \textbf{推论1（Burnside）}：设不等价不可约表示维数为$n_r$，则
  \[
  \sum_r n_r^{2}=|G|.
  \]
  \item \textbf{推论2（正则表示的约化）}：正则表示可完全约化为所有不等价不可约表示的直和，且每个不可约表示出现次数等于其维数：
  \[
  \mathrm{Reg}\ \cong\ \bigoplus_{r} n_r\,A^{(r)}.
  \]
\end{itemize}
\noindent\textcolor{red}{\rule{\linewidth}{1pt}}
\section*{特征标（Character）}
\begin{itemize}
  \item \textbf{定义}：表示$A$的特征标（character）为$\chi^{(A)}(g)\equiv \mathrm{tr}\,A(g)$，是一个群函数。
  \item \textbf{等价不变}：若$A'(g)=X A(g)X^{-1}$，则$\chi^{(A')}(g)=\chi^{(A)}(g)$。
  \item \textbf{类函数}：若$h=gfg^{-1}$与$f$共轭，则$\chi^{(A)}(h)=\chi^{(A)}(f)$（同一共轭类取值相同）。
\end{itemize}

\section*{特征标的第一正交关系}
\begin{itemize}
  \item 设$G$为$n$阶有限群，$\{\chi^{(r)}\}$为不等价不可约表示的特征标，则
  \[
  \left\langle \chi^{(p)}\mid \chi^{(r)} \right\rangle
  \equiv \frac{1}{n}\sum_{g\in G}\big(\chi^{(p)}(g)\big)^{*}\chi^{(r)}(g)=\delta_{pr}.
  \]
  \item 按共轭类求和：若$G$有共轭类$K_i$，大小$|K_i|=n_i$，则
  \[
  \frac{1}{n}\sum_{i} n_i\,\big(\chi^{(p)}(K_i)\big)^{*}\chi^{(r)}(K_i)=\delta_{pr}.
  \]
  \item 特别地，$\langle\chi^{(r)}|\chi^{(r)}\rangle=1$（不可约特征标的“范数”）。
\end{itemize}

\section*{用特征标判断可约性与重数}
\begin{itemize}
  \item 若表示$A$完全可约分解为$A\simeq \bigoplus_{p} m_p A^{(p)}$，则
  \[
  m_p=\left\langle \chi^{(p)}\mid \chi^{(A)} \right\rangle.
  \]
  \item \textbf{判别}：$\langle\chi^{(A)}|\chi^{(A)}\rangle=\sum_p m_p^2$；因此
  \[
  \langle\chi^{(A)}|\chi^{(A)}\rangle=1 \Rightarrow A\ \text{不可约},\qquad
  \langle\chi^{(A)}|\chi^{(A)}\rangle>1 \Rightarrow A\ \text{可约}.
  \]
\end{itemize}

\section*{类函数与特征标完备性}
\begin{itemize}
  \item \textbf{类函数}：群函数$f$若对任意共轭$h=gfg^{-1}$满足$f(h)=f(f)$，则称$f$为类函数；全体类函数构成群函数空间的子空间。
  \item \textbf{线性独立}：有限群不等价不可约表示的特征标$\{\chi^{(r)}\}$是线性独立的类函数。
  \item \textbf{特征标完备性定理}：有限群不等价不可约表示的特征标在类函数空间中完备；任意类函数可展开为$f=\sum_r a^{(r)}\chi^{(r)}$。
  \item \textbf{推论}：有限群不等价不可约表示的个数$=$共轭类个数。
\end{itemize}

\section*{特征标的第二正交关系}
\begin{itemize}
  \item 设$G$为$n$阶有限群，有共轭类$K_i$（大小$|K_i|=n_i$），则
  \[
  \sum_{r}\big(\chi^{(r)}(K_i)\big)^{*}\chi^{(r)}(K_j)=\frac{n}{n_i}\,\delta_{ij}.
  \]
\end{itemize}
\noindent\textcolor{red}{\rule{\linewidth}{1pt}}
\section*{特征标表（Character Table）}
\begin{itemize}
  \item 将有限群$G$的全部不等价不可约表示的特征标$\chi^{(r)}$按\textbf{行}排列、共轭类$K_i$按\textbf{列}排列，得到的表称为$G$的特征标表。
  \item \textbf{第一行}通常为一维恒等表示，整行全为$1$；\textbf{第一列}对应单位元类$\{e\}$，其值为表示维数：$\chi^{(r)}(e)=d_r$。
  \item 行/列正交：行满足第一正交关系，列满足第二正交关系（对应“行正交/列正交”）。
\end{itemize}
\begin{tabular}{c|cccc}
  & $n_1 K_1$ & $n_2 K_2$ & $\cdots$ & $n_q K_q$ \\
 \hline
 $A^{(1)}$ & $\chi^{(1)}(K_1)$ & $\chi^{(1)}(K_2)$ & $\cdots$ & $\chi^{(1)}(K_q)$ \\
 $A^{(2)}$ & $\chi^{(2)}(K_1)$ & $\chi^{(2)}(K_2)$ & $\cdots$ & $\chi^{(2)}(K_q)$ \\
 $\vdots$ & $\vdots$ & $\vdots$ & $\ddots$ & $\vdots$ \\
 $A^{(q)}$ & $\chi^{(q)}(K_1)$ & $\chi^{(q)}(K_2)$ & $\cdots$ & $\chi^{(q)}(K_q)$ \\
 \end{tabular}
\section*{阿贝尔群的不可约表示}
\begin{itemize}
  \item 阿贝尔群每个元素自成共轭类，故不等价不可约表示数$=|G|$。
  \item 由$\sum_r d_r^2=|G|$可知阿贝尔群的全部不可约表示均为$1$维（$d_r=1$）。
\end{itemize}
\noindent\textcolor{red}{\rule{\linewidth}{1pt}}
%======================================
% 置换群专题
%======================================
\section*{置换的定义与表示}
\begin{itemize}
  \item \textbf{定义}：将$n$个物体（标记为$1,2,\dots,n$）的一个置换操作记为：
  \[
  P=\begin{pmatrix} 1 & 2 & \cdots & n \\ m_1 & m_2 & \cdots & m_n \end{pmatrix}
  \]
  其中$m_1,\dots,m_n$是$1,\dots,n$的某种排列。表示$1\to m_1,\ 2\to m_2,\ \dots,\ n\to m_n$。
  \item \textbf{列无关性}：置换$P$中各列的相对位置变化不会改变该映射。
  \item \textbf{逆置换}：$P$的逆$P^{-1}$为
  \[
  P^{-1}=\begin{pmatrix} m_1 & m_2 & \cdots & m_n \\ 1 & 2 & \cdots & n \end{pmatrix}
  \]
  即交换上下行后重新排列列顺序。
\end{itemize}

\section*{置换的乘法与$S_n$群}
\begin{itemize}
  \item \textbf{恒等置换}：$P_0=\begin{pmatrix} 1 & \cdots & n \\ 1 & \cdots & n \end{pmatrix}$。
  \item \textbf{乘积定义}：两个置换的乘积$P_1 P_2$是\textbf{先操作右边}的置换$P_2$，\textbf{再操作左边}的置换$P_1$。
  \item \textbf{非交换性}：一般地，$P_1 P_2 \neq P_2 P_1$。
  \item \textbf{对称群$S_n$}：$n$个物体的全体置换在上述乘法定义下构成一个群，称为$n$阶对称群，记为$S_n$。
  \item \textbf{阶数}：$S_n$的阶（元素个数）为$n!$。
\end{itemize}

\section*{轮换 (Cycle)}
\begin{itemize}
  \item \textbf{定义}：$m$个物体$a_1, a_2, \dots, a_m$的如下置换称为一个\textbf{轮换}：
  \[
  \begin{pmatrix} a_1 & a_2 & \cdots & a_m \\ a_2 & a_3 & \cdots & a_1 \end{pmatrix} \equiv (a_1 a_2 \cdots a_m)
  \]
  \item \textbf{映射关系}：定义了一个循环映射 $a_1 \to a_2 \to a_3 \to \cdots \to a_m \to a_1$。
  \item \textbf{长度}：$m$称为这个轮换的长度。
  \item \textbf{轮换的逆}：$(a_1 a_2 \cdots a_m)^{-1} = (a_m a_{m-1} \cdots a_1)$。
\end{itemize}

\section*{轮换的分解性质}
\begin{itemize}
  \item \textbf{轮换分解定理}：{\color{red}任意置换都可\textbf{唯一}地表达为\textbf{不含公共数字}（不相交）的轮换的乘积。}
  \item \textbf{轮换结构}：该表达方式称为置换的轮换结构，由轮换的\textbf{长度}和\textbf{个数}决定。
  \item \textbf{可交换性}：两个不相交的轮换是\textbf{可交换}的。
  \[ \text{例：} (1\,4)(2\,3\,5) = (2\,3\,5)(1\,4) \]
  \item \textbf{矩阵转轮换示例}：
  \[
  \begin{pmatrix} 1 & 2 & 3 & 4 & 5 \\ 4 & 3 & 5 & 1 & 2 \end{pmatrix} = (1\,4)(2\,3\,5)
  \]
\end{itemize}

\section*{对换与生成元}
\begin{itemize}
  \item \textbf{对换 (Transposition)}：长度为2的轮换$(i\,j)$。
  \item \textbf{轮换拆解}：{\color{red}长度为$m$的轮换可写成$m-1$个对换的乘积}（形式不唯一）：
  \begin{align*}
  (a_1 a_2 \cdots a_m) &= (a_1 a_m)(a_1 a_{m-1})\cdots(a_1 a_2) \\
                       &= (a_1 a_2)(a_2 a_3)\cdots(a_{m-1} a_m)
  \end{align*}
  \item \textbf{引理1}：$(a_1, a_2 \cdots a_m) = (a_1, a_m)(a_1, a_2 \cdots a_{m-1})$
  \item \textbf{引理2}：$(a_1, a_2 \cdots a_m) = (a_1, a_2)(a_2, a_3 \cdots a_m)$
  \item \textbf{邻换 (Adjacent Transposition)}：相邻数字的对换$(k, k+1)$。
  \item \textbf{任意对换的分解}：{\color{red}任意对换都可写成邻换的乘积}。
  \[ (1\,3) = (2\,3)(1\,2)(2\,3) \]
  \item \textbf{群的生成}：
  \begin{enumerate}
      \item 任意置换可由对换生成；
      \item $S_n$群可由$n-1$个邻换生成：$(1,2), (2,3), \dots, (n-1, n)$。
  \end{enumerate}
\end{itemize}

\section*{置换的奇偶性}
\begin{itemize}
  \item \textbf{定义}：一个置换可表达为\textbf{奇数}个对换的乘积，称为\textbf{奇置换}；若为\textbf{偶数}个，称为\textbf{偶置换}。
  \item \textbf{唯一性 (性质4)}：{\color{red}尽管对换分解形式不唯一，但分解出的对换个数的奇偶性是唯一确定}的。
  \item \textbf{逆元性质 (性质5)}：{\color{red}任意置换$P$和它的逆$P^{-1}$具有相同的奇偶性}。
  \item \textbf{关系}：
  \[ \text{偶} \times \text{偶} = \text{偶}, \quad \text{奇} \times \text{奇} = \text{偶}, \quad \text{奇} \times \text{偶} = \text{奇} \]
\end{itemize}

\section*{实例：$S_3$ 群结构}
\begin{itemize}
  \item \textbf{元素}：$S_3 = \{e, (12), (23), (13), (123), (132)\}$，共$3!=6$个。
  \item \textbf{对应关系}（与$D_3$同构）：
  \begin{itemize}
      \item 偶置换 (旋转)：$e$ (单位元), $d=(123)$, $f=(132)$
      \item 奇置换 (反射)：$a=(12)$, $b=(23)$, $c=(13)$
  \end{itemize}
  \item \textbf{乘法示例}（右乘优先）：
  \[ (1\,2\,3)(1\,2) = (1\,3) \quad (\text{即 } da=c) \]
  \[ (1\,2\,3)(2\,3) = (1\,2) \quad (\text{即 } db=a) \]
  \item \textbf{生成元}：$S_3$可由两个邻换$(1\,2)$和$(2\,3)$生成。
\end{itemize}
\section*{置换的共轭}
\begin{itemize}
  \item \textbf{共轭操作}：用$Q$对$P$作共轭变换 $QPQ^{-1}$。
  \item \textbf{几何意义}：若$P=\begin{pmatrix} x \\ P(x) \end{pmatrix}$，则$QPQ^{-1}$ 相当于对$P$的上下两行数字同时作$Q$置换（即对元素进行重命名）。
  \item \textbf{性质6}：{\color{red}任意两个对换都是互相共轭的}。
  \item \textbf{性质7}：{\color{red}互相共轭的两个置换具有相同的轮换结构}（即轮换的长度和个数分布相同）。
  \item \textbf{逆命题}：具有相同轮换结构的置换属于同一个共轭类。
  \item \textbf{结论}：$S_n$中，\textbf{共轭类}由\textbf{轮换结构}唯一确定。
\end{itemize}

\section*{轮换结构与分划}
\begin{itemize}
  \item \textbf{标记 $(\nu)$}：$S_n$群的轮换结构可记为 $(\nu) = (1^{\nu_1}, 2^{\nu_2}, \cdots, n^{\nu_n})$。
  \item \textbf{含义}：$\nu_k$ 表示长度为 $k$ 的\textbf{独立}轮换的个数。
  \item \textbf{约束}：非负整数 $\nu_k$ 满足 $\sum_{k=1}^n k \cdot \nu_k = n$。
  \item \textbf{分划 (Partition)}：也可描述为 $n$ 的分划 $[\lambda] = [\lambda_1, \lambda_2, \cdots, \lambda_m]$，满足 $\sum \lambda_i = n$ 且 $\lambda_1 \ge \lambda_2 \ge \dots \ge 0$。
\end{itemize}

\section*{共轭类元素个数公式}
\begin{itemize}
  \item 对于轮换结构 $(\nu)=(1^{\nu_1}, 2^{\nu_2}, \cdots, n^{\nu_n})$，该共轭类包含的元素数目 $\rho_{(\nu)}$ 为：
  \[
  \rho_{(\nu)} = \frac{n!}{1^{\nu_1}\nu_1! \, 2^{\nu_2}\nu_2! \, \cdots \, n^{\nu_n}\nu_n!} = \frac{n!}{\prod_{k=1}^n k^{\nu_k} \nu_k!}
  \]
  \item \textbf{例子 ($S_6$)}：
  \begin{itemize}
      \item 结构 $(******)$ 即 $(6^1)$：$\frac{6!}{6^1 1!} = 120$。
      \item 结构 $(**)(**)(**)$ 即 $(2^3)$：$\frac{6!}{2^3 3!} = \frac{720}{8 \times 6} = 15$。
  \end{itemize}
\end{itemize}

\section*{实例：$S_3$ 群的共轭类}
$S_3$ 共有3个共轭类，对应3种分划：
\begin{center}
\begin{tabular}{|c|c|c|c|}
\hline
\textbf{轮换结构} & \textbf{标记} $(\nu)$ & \textbf{分划} $[\lambda]$ & \textbf{元素数} \\
\hline
$(*)(*)(*)$ & $(1^3)$ & $[3]$ & $\frac{3!}{1^3 3!} = 1$ \\
\hline
$(*)(**)$ & $(1^1, 2^1)$ & $[2,1]$ & $\frac{3!}{1^1 1! \, 2^1 1!} = 3$ \\
\hline
$(***)$ & $(3^1)$ & $[1,1,1]$ & $\frac{3!}{3^1 1!} = 2$ \\
\hline
\end{tabular}
\end{center}

\section*{实例：$S_4$ 群的共轭类}
$S_4$ 群的共轭类由 4 的分划 $[\lambda]$ 决定，共有 5 种情况。
\begin{enumerate}
    \item $[\lambda]=[4]$：对应轮换结构 $(1^4)$，即 $(\ast)(\ast)(\ast)(\ast)$（单位元）。\\
    元素数目：$1$。
    \item $[\lambda]=[3,1]$：对应轮换结构 $(1^2, 2^1)$，即 $(\ast)(\ast)(\ast\ast)$（对换）。\\
    元素数目：$\frac{4!}{1^2 \cdot 2! \cdot 2^1 \cdot 1!} = 6$。
    \item $[\lambda]=[2,2]$：对应轮换结构 $(2^2)$，即 $(\ast\ast)(\ast\ast)$。\\
    元素数目：$\frac{4!}{2^2 \cdot 2!} = 3$。
    \item $[\lambda]=[2,1,1]$：对应轮换结构 $(1^1, 3^1)$，即 $(\ast)(\ast\ast\ast)$。\\
    元素数目：$\frac{4!}{1^1 \cdot 1! \cdot 3^1 \cdot 1!} = 8$。
    \item $[\lambda]=[1,1,1,1]$：对应轮换结构 $(4^1)$，即 $(\ast\ast\ast\ast)$。\\
    元素数目：$\frac{4!}{4^1 \cdot 1!} = 6$。
\end{enumerate}
\textbf{验证}：$1+6+3+8+6 = 24 = 4!$。

\section*{杨图 (Young Diagram)}
\begin{itemize}
    \item \textbf{定义}：自然数 $n$ 的一个分划 $[\lambda] = [\lambda_1, \lambda_2, \dots, \lambda_m]$ 可用杨图刻画。它由 $n$ 个方格组成，第 1 行有 $\lambda_1$ 个，第 2 行 $\lambda_2$ 个，以此类推。
    \item \textbf{$S_3$ 的杨图}：
    \begin{itemize}
        \item $[3]$: \ydiagram{3}
        \item $[2,1]$: \ydiagram{2,1}
        \item $[1,1,1]$: \ydiagram{1,1,1}
    \end{itemize}
    \item \textbf{$S_4$ 的杨图}：
    \begin{multicols}{3}
        $[4]$: \ydiagram{4} \\
        $[3,1]$: \ydiagram{3,1} \\
        $[2,2]$: \ydiagram{2,2}
        
        \columnbreak
        
        $[2,1,1]$: \ydiagram{2,1,1} \\
        $[1,1,1,1]$: \ydiagram{1,1,1,1}
    \end{multicols}
\end{itemize}

\section*{对偶杨图}
\begin{itemize}
    \item \textbf{定义}：将杨图 $[\lambda]$ 的行与列互换，得到杨图 $[\tilde{\lambda}]$，称两者互为对偶。
    \item \textbf{例子}：
    \begin{itemize}
        \item $S_3$ 中，$[3]$ 与 $[1,1,1]$ 对偶；$[2,1]$ 是自对偶。
        \item $S_4$ 中，$[4]$ 与 $[1,1,1,1]$ 对偶；$[3,1]$ 与 $[2,1,1]$ 对偶；$[2,2]$ 自对偶。
    \end{itemize}
\end{itemize}

\section*{杨盘 (Young Tableau)}
\begin{itemize}
    \item \textbf{杨盘}：将数字 $1, 2, \dots, n$ 不重复地填入杨图 $[\lambda]$ 的 $n$ 个方格中。
    \item \textbf{标准杨盘 (SYT)}：
    \begin{enumerate}
        \item 每一行从左到右递增；
        \item 每一列从上到下递增。
    \end{enumerate}
    记为 $T_{\gamma}^{[\lambda]}$，其中 $\gamma$ 标记不同的标准杨盘。
\end{itemize}

\section*{标准杨盘实例}
\textbf{1. $S_3$ 群 (共4个)}
\begin{itemize}
    \item $T^{[3]}$ (1个): \begin{ytableau} 1 & 2 & 3 \end{ytableau}
    \item $T^{[2,1]}$ (2个): 
    \begin{ytableau} 1 & 2 \\ 3 \end{ytableau} \quad
    \begin{ytableau} 1 & 3 \\ 2 \end{ytableau}
    \item $T^{[1,1,1]}$ (1个): \begin{ytableau} 1 \\ 2 \\ 3 \end{ytableau}
\end{itemize}
\emph{注：对偶杨图的标准杨盘数目相等。}

\textbf{2. $S_4$ 群 (共10个)}
\begin{itemize}
    \item $T^{[4]}$: \begin{ytableau} 1 & 2 & 3 & 4 \end{ytableau}
    \item $T^{[3,1]}$ (3个): 
    \tiny
    \begin{ytableau} 1 & 2 & 3 \\ 4 \end{ytableau} \,
    \begin{ytableau} 1 & 2 & 4 \\ 3 \end{ytableau} \,
    \begin{ytableau} 1 & 3 & 4 \\ 2 \end{ytableau}
    \normalsize
    \item $T^{[2,2]}$ (2个):
    \tiny
    \begin{ytableau} 1 & 2 \\ 3 & 4 \end{ytableau} \,
    \begin{ytableau} 1 & 3 \\ 2 & 4 \end{ytableau}
    \normalsize
    \item $T^{[2,1,1]}$ (3个):
    \tiny
    \begin{ytableau} 1 & 2 \\ 3 \\ 4 \end{ytableau} \,
    \begin{ytableau} 1 & 3 \\ 2 \\ 4 \end{ytableau} \,
    \begin{ytableau} 1 & 4 \\ 2 \\ 3 \end{ytableau}
    \normalsize
    \item $T^{[1,1,1,1]}$: 
    \tiny \begin{ytableau} 1 \\ 2 \\ 3 \\ 4 \end{ytableau} \normalsize
\end{itemize}

\section*{不可约表示维数与钩长公式}
\begin{itemize}
    \item \textbf{定理}：$S_n$ 群杨图 $[\lambda]$ 对应的不可约表示的维数 $d_{[\lambda]}$，等于该杨图的标准杨盘数目 $f^{[\lambda]}$。
    \item \textbf{钩长 (Hook Length)}：杨图第 $i$ 行第 $j$ 列方格的钩长 $g_{ij}$ 定义为：
    \[ g_{ij} = (\text{右边方格数}) + (\text{下边方格数}) + 1 \]
    \item \textbf{钩长公式}：
    \[ f^{[\lambda]} = \frac{n!}{\prod_{(i,j)} g_{ij}} \]
\end{itemize}

\section*{实例：$S_3$ 不可约表示维数}
$S_3$ 阶数为 $3!=6$。计算各杨图钩长（填入格中）及维数：
\begin{multicols}{3}
    \textbf{1. [3]} \\
    \begin{ytableau} 3 & 2 & 1 \end{ytableau} \\
    $f^{[3]} = \frac{6}{3\cdot 2\cdot 1} = 1$

    \columnbreak

    \textbf{2. [2,1]} \\
    \begin{ytableau} 3 & 1 \\ 1 \end{ytableau} \\
    $f^{[2,1]} = \frac{6}{3\cdot 1\cdot 1} = 2$

    \columnbreak

    \textbf{3. [1,1,1]} \\
    \begin{ytableau} 3 \\ 2 \\ 1 \end{ytableau} \\
    $f^{[1^3]} = \frac{6}{6} = 1$
\end{multicols}
\textbf{验证}：$1^2 + 2^2 + 1^2 = 6 = |S_3|$。

% \section*{实例：$S_4$ 不可约表示维数}
% $S_4$ 阶数为 $4!=24$。计算如下：
% \begin{itemize}
%     \item \textbf{[4]}：钩长 \tiny \begin{ytableau} 4 & 3 & 2 & 1 \end{ytableau} \normalsize $\Rightarrow f = \frac{24}{24} = 1$。
    
%     \item \textbf{[3,1]}：钩长 \tiny \begin{ytableau} 4 & 2 & 1 \\ 1 \end{ytableau} \normalsize $\Rightarrow f = \frac{24}{4\cdot 2\cdot 1\cdot 1} = 3$。
    
%     \item \textbf{[2,2]}：钩长 \tiny \begin{ytableau} 3 & 2 \\ 2 & 1 \end{ytableau} \normalsize $\Rightarrow f = \frac{24}{3\cdot 2\cdot 2\cdot 1} = 2$。
    
%     \item \textbf{[2,1,1]}：钩长 \tiny \begin{ytableau} 4 & 1 \\ 2 \\ 1 \end{ytableau} \normalsize $\Rightarrow f = \frac{24}{4\cdot 1\cdot 2\cdot 1} = 3$。
    
%     \item \textbf{[1,1,1,1]}：钩长 \tiny \begin{ytableau} 4 \\ 3 \\ 2 \\ 1 \end{ytableau} \normalsize $\Rightarrow f = \frac{24}{24} = 1$。
% \end{itemize}
% \textbf{验证}：$1^2 + 3^2 + 2^2 + 3^2 + 1^2 = 1+9+4+9+1 = 24 = |S_4|$。
%======================================
% 杨算符与基函数构造
%======================================
\section*{杨盘之间的联系}
\begin{itemize}
    \item \textbf{置换联系}：同一杨图 $[\lambda]$ 的两个杨盘 $T_\gamma$ 和 $T_s$ 可通过置换 $\sigma_{\gamma s} \in S_n$ 联系：$T_\gamma = \sigma_{\gamma s} T_s$。
    \item \textbf{对换示例}：对于 $T_1^{[2,1]} =$ \tiny\begin{ytableau}1 & 2 \\ 3\end{ytableau} 
    和 $T_2^{[2,1]} =$ \tiny\begin{ytableau}1 & 3 \\ 2\end{ytableau}，有 $T_1 = (23) T_2$。
\end{itemize}

\section*{行置换与列置换集合}
给定一个杨盘 $T$，定义：
\begin{itemize}
    \item \textbf{行置换集合 $R(T)$}：保持每一行内元素集合不变的所有置换（即行内元素的任意交换）。
    \item \textbf{列置换集合 $C(T)$}：保持每一列内元素集合不变的所有置换。
    \item \textbf{示例}：$T =$ \tiny\begin{ytableau}1 & 2 \\ 3 & 4\end{ytableau}，
    则 $R(T)=\{e, (12), (34), (12)(34)\}$，$C(T)=\{e, (13), (24), (13)(24)\}$。
\end{itemize}

\section*{杨算符 (Young Operator)}
\begin{itemize}
    \item \textbf{对称化算符 $P(T)$}：行置换之和。
    $P(T) = \sum_{p \in R(T)} p$
    \item \textbf{反对称化算符 $Q(T)$}：列置换的加权和（偶置换系数+1，奇置换-1）。
    $Q(T) = \sum_{q \in C(T)} \delta_q q = \sum_{q \in C(T)} (-1)^{\text{sgn}(q)} q$
    \item \textbf{杨算符 $E(T)$}：先对称化后反对称化。
    $E(T) = P(T)Q(T)$
\end{itemize}

\section*{实例：$S_3$ 的不可约表示构造}
利用杨算符作用于初始函数 $\psi(1,2,3)$ 可生成不可约表示的基函数。

\textbf{1. 表示 [3] (全对称)}
\begin{itemize}
    \item 杨盘 $T^{[3]} =$ \tiny\begin{ytableau}1 & 2 & 3\end{ytableau}。
    \item $R(T) = S_3,\ C(T) = \{e\}$。
    \item $E(T) = \sum_{p \in S_3} p$。
    \item 基函数 $\psi_s = E(T)\psi$ 为\textbf{全对称函数}。对应 1 维恒等表示。
\end{itemize}

\textbf{2. 表示 [1,1,1] (全反对称)}
\begin{itemize}
    \item 杨盘 $T^{[1^3]} =$ \tiny\begin{ytableau}1 \\ 2 \\ 3\end{ytableau}。
    \item $R(T) = \{e\},\ C(T) = S_3$。
    \item $E(T) = \sum_{q \in S_3} \delta_q q$。
    \item 基函数 $\psi_a = E(T)\psi$ 为\textbf{全反对称函数}。对应 1 维交错表示（偶+1，奇-1）。
\end{itemize}

\textbf{3. 表示 [2,1] (二维表示)}
\begin{itemize}
    \item 杨盘 $T_1 =$ \tiny\begin{ytableau}1 & 2 \\ 3\end{ytableau}，$T_2 =$ \tiny\begin{ytableau}1 & 3 \\ 2\end{ytableau}。
    \item 对于 $T_1$：$P(T_1)=1+(12),\ Q(T_1)=1-(13)$。
    \[ E(T_1) = [1+(12)][1-(13)] = 1+(12)-(13)-(123) \]
    基函数 $\Psi_1 = E(T_1)\psi(1,2,3)$。
    \item 对于 $T_2$：$P(T_2)=1+(13),\ Q(T_2)=1-(12)$。
    \[ E(T_2) = [1+(13)][1-(12)] = 1+(13)-(12)-(132) \]
    基函数 $\Psi_2 = E(T_2)\psi_2$，其中 $\psi_2=(23)\psi(1,2,3)$。
    \item $\{\Psi_1, \Psi_2\}$ 构成 $S_3$ 群 2 维不可约表示的一组基。
\end{itemize}
%======================================
% S3群二维表示的矩阵推导与幺正化
%======================================
\section*{二维不可约表示 [2,1] 的矩阵形式}
利用前述基函数 $\Psi_1, \Psi_2$，考察群元素的作用：
\begin{itemize}
    \item \textbf{生成元 (12)}：
    \[
    \begin{cases} (12)\Psi_1 = \Psi_1 \\ (12)\Psi_2 = -\Psi_1 - \Psi_2 \end{cases}
    \implies A(12) = \begin{pmatrix} 1 & -1 \\ 0 & -1 \end{pmatrix}
    \]
    \item \textbf{生成元 (13)}：
    \[
    \begin{cases} (13)\Psi_1 = -\Psi_1 - \Psi_2 \\ (13)\Psi_2 = \Psi_2 \end{cases}
    \implies A(13) = \begin{pmatrix} -1 & 0 \\ -1 & 1 \end{pmatrix}
    \]
    \item \textbf{其他元素}（通过矩阵乘法生成）：
    \begin{align*}
    A(123) &= A(13)A(12) = \begin{pmatrix} -1 & 1 \\ -1 & 0 \end{pmatrix} \\
    A(321) &= A(12)A(13) = \begin{pmatrix} 0 & -1 \\ 1 & -1 \end{pmatrix} \\
    A(23) &= A(12)A(13)A(12) = \begin{pmatrix} 0 & 1 \\ 1 & 0 \end{pmatrix}
    \end{align*}
\end{itemize}

\section*{矩阵汇总表}
$S_3$ 群 2 维不可约表示 [2,1] 的所有矩阵：
\begin{center}
\renewcommand{\arraystretch}{1.5}
\begin{tabular}{|c|c|c|}
    \hline
    $e$ & $(12)$ & $(13)$ \\
    $\begin{pmatrix} 1 & 0 \\ 0 & 1 \end{pmatrix}$ & $\begin{pmatrix} 1 & -1 \\ 0 & -1 \end{pmatrix}$ & $\begin{pmatrix} -1 & 0 \\ -1 & 1 \end{pmatrix}$ \\
    \hline
    $(23)$ & $(123)$ & $(321)$ \\
    $\begin{pmatrix} 0 & 1 \\ 1 & 0 \end{pmatrix}$ & $\begin{pmatrix} -1 & 1 \\ -1 & 0 \end{pmatrix}$ & $\begin{pmatrix} 0 & -1 \\ 1 & -1 \end{pmatrix}$ \\
    \hline
\end{tabular}
\end{center}

\section*{不可约性验证 (特征标法)}
计算上述矩阵的迹（特征标 $\chi$）：
\begin{itemize}
    \item $\chi(e) = 1+1 = 2$
    \item $\chi((123)) = -1+0 = -1$ （同类元素 $(321)$ 相同）
    \item $\chi((12)) = 1-1 = 0$ （同类元素 $(13), (23)$ 相同）
\end{itemize}
计算特征标模长平方：
\[
\langle \chi | \chi \rangle = \frac{1}{6} \left[ 1 \cdot (2)^2 + 2 \cdot (-1)^2 + 3 \cdot (0)^2 \right] = \frac{1}{6}(4+2) = 1
\]
因为 $\langle \chi | \chi \rangle = 1$，故该表示是\textbf{不可约}的。

\section*{表示的幺正化}
上述表示\textbf{不是}幺正表示。例如验证 $A(12)$：
\[ A(12) A(12)^\dagger = \begin{pmatrix} 1 & -1 \\ 0 & -1 \end{pmatrix} \begin{pmatrix} 1 & 0 \\ -1 & -1 \end{pmatrix} = \begin{pmatrix} 2 & 1 \\ 1 & 1 \end{pmatrix} \neq I \]
\textbf{幺正化步骤}：
\begin{enumerate}
    \item \textbf{构造厄米矩阵 $S$}：
    \begin{align*}
    S &= \sum_{g \in G} [A(g)]^\dagger A(g) \\
      &= I + \begin{pmatrix} 1 & 0 \\ -1 & -1 \end{pmatrix}\begin{pmatrix} 1 & -1 \\ 0 & -1 \end{pmatrix} + \dots (\text{共6项}) \\
      &= \begin{pmatrix} 8 & -4 \\ -4 & 8 \end{pmatrix}
    \end{align*}
    \item \textbf{求 $S$ 的平方根 $X = S^{1/2}$}：
    对角化 $S$ 得本征值 $12, 4$。计算得：
    \[ X = \begin{pmatrix} \frac{1+\sqrt{3}}{2} & \frac{1-\sqrt{3}}{2} \\ \frac{1-\sqrt{3}}{2} & \frac{1+\sqrt{3}}{2} \end{pmatrix} \]
    （此处略去中间对角化过程，仅展示结果形式）
    \item \textbf{相似变换}：
    新的表示 $A'(g) = X A(g) X^{-1}$ 即为幺正表示。
\end{enumerate}
%======================================
% 通用规则：从杨图读出正交表示矩阵
%======================================
\section*{通用规则：直接读出矩阵元}
$S_n$ 群生成元 $(k, k+1)$ 的矩阵元可由标准杨盘（SYT）直接定出。设杨图 $[\lambda]$ 有标准杨盘 $T_1, T_2, \dots$。
\begin{enumerate}
    \item \textbf{对角元}：
    \begin{itemize}
        \item 若 $k$ 和 $k+1$ 在 $T_r$ 中位于\textbf{同一行}：$(k, k+1)_{rr} = 1$。
        \item 若 $k$ 和 $k+1$ 在 $T_r$ 中位于\textbf{同一列}：$(k, k+1)_{rr} = -1$。
    \end{itemize}
    \item \textbf{非对角元}：
    若 $k, k+1$ 既不同行也不同列，设 $(k, k+1)$ 将 $T_r$ 变为 $T_s$（即交换数字位置后仍为标准杨盘），且在 $T_r$ 中 $k$ 到 $k+1$ 的\textbf{轴距离}为 $d$（令 $\rho = 1/d$），则：
    \[
    \begin{pmatrix} (k, k+1)_{rr} & (k, k+1)_{rs} \\ (k, k+1)_{sr} & (k, k+1)_{ss} \end{pmatrix} = 
    \begin{pmatrix} -\rho & \sqrt{1-\rho^2} \\ \sqrt{1-\rho^2} & \rho \end{pmatrix}
    \]
\end{enumerate}

\textbf{轴距离 (Axial Distance) 定义}：
从数字 $a$ 开始数格子到 $b$：
\begin{itemize}
    \item 向\textbf{左}或向\textbf{下}数 1 格：$+1$
    \item 向\textbf{右}或向\textbf{上}数 1 格：$-1$
\end{itemize}
最后得到的步数之和即为 $d = 1/\rho$。
\begin{itemize}
    \item 例：\tiny\begin{ytableau} \none & a \\ b \end{ytableau} ($a \to b$: 左1下1)，$d=2 \Rightarrow \rho=1/2$。
    \item 例：\tiny\begin{ytableau} \none & b \\ a \end{ytableau} ($a \to b$: 右1上1)，$d=-2 \Rightarrow \rho=-1/2$。
\end{itemize}

\section*{应用：$S_3$ 群 [2,1] 表示 (重算)}
标准杨盘：$T_1 =$ \tiny\begin{ytableau} 1 & 2 \\ 3 \end{ytableau}, \quad $T_2 =$ \tiny\begin{ytableau} 1 & 3 \\ 2 \end{ytableau}。
\begin{itemize}
    \item \textbf{(12)}：
    \begin{itemize}
        \item $T_1$: 1,2 同行 $\to 1$。
        \item $T_2$: 1,2 同列 $\to -1$。
        \item 矩阵：$\begin{pmatrix} 1 & 0 \\ 0 & -1 \end{pmatrix}$。
    \end{itemize}
    \item \textbf{(23)}：
    \begin{itemize}
        \item $T_1$ 中 2,3 不同行不同列。$2(1,2) \to 3(2,1)$。
        \item 路径：左1 ($+1$)，下1 ($+1$)。距离 $d=2 \Rightarrow \rho=1/2$。
        \item 矩阵：$\begin{pmatrix} -1/2 & \sqrt{3}/2 \\ \sqrt{3}/2 & 1/2 \end{pmatrix}$。
    \end{itemize}
\end{itemize}

\section*{应用：$S_4$ 群 [3,1] 表示 (3维)}
标准杨盘：
$T_1=$ \tiny\begin{ytableau}1&2&3\\4\end{ytableau}, $T_2=$ \tiny\begin{ytableau}1&2&4\\3\end{ytableau}, $T_3=$ \tiny\begin{ytableau}1&3&4\\2\end{ytableau}。
\begin{itemize}
    \item \textbf{(12)}：
    $T_1(1), T_2(1), T_3(-1)$ $\Rightarrow \text{diag}(1, 1, -1)$。
    \item \textbf{(23)}：
    \begin{itemize}
        \item $T_1$: 2,3 同行 $\to 1$。
        \item $T_2 \leftrightarrow T_3$: $2(1,2) \to 3(2,1)$，距离 $d=2 \Rightarrow \rho=1/2$。
        \item 矩阵：$\begin{pmatrix} 1 & 0 & 0 \\ 0 & -1/2 & \sqrt{3}/2 \\ 0 & \sqrt{3}/2 & 1/2 \end{pmatrix}$。
    \end{itemize}
    \item \textbf{(34)}：
    \begin{itemize}
        \item $T_3$: 3,4 同行 $\to 1$。
        \item $T_1 \leftrightarrow T_2$: $3(1,3) \to 4(2,1)$。
        \item 路径：左1(+1), Left(+1), Down(+1). Wait, 3 is at (1,3), 4 is at (2,1).
        \item Path: Left(to 1,2)+1, Left(to 1,1)+1, Down(to 2,1)+1. Sum=3. $\rho=1/3$。
        \item 矩阵：$\begin{pmatrix} -1/3 & \sqrt{8}/3 & 0 \\ \sqrt{8}/3 & 1/3 & 0 \\ 0 & 0 & 1 \end{pmatrix}$。
    \end{itemize}
\end{itemize}

\section*{应用：$S_4$ 群 [2,2] 表示 (2维)}
$T_1=$ \tiny\begin{ytableau}1&2\\3&4\end{ytableau}, $T_2=$ \tiny\begin{ytableau}1&3\\2&4\end{ytableau}。
\begin{itemize}
    \item \textbf{(12)}: $T_1(1), T_2(-1) \Rightarrow \text{diag}(1, -1)$。
    \item \textbf{(34)}: $T_1(1), T_2(-1) \Rightarrow \text{diag}(1, -1)$。
    \item \textbf{(23)}: $T_1 \leftrightarrow T_2$. $2(1,2) \to 3(2,1)$. $d=2 \Rightarrow \rho=1/2$.
    \item 矩阵：$\begin{pmatrix} -1/2 & \sqrt{3}/2 \\ \sqrt{3}/2 & 1/2 \end{pmatrix}$。
\end{itemize}

\section*{应用：$S_4$ 群 [2,1,1] 表示 (3维)}
$T_1=$ \tiny\begin{ytableau}1&2\\3\\4\end{ytableau}, $T_2=$ \tiny\begin{ytableau}1&3\\2\\4\end{ytableau}, $T_3=$ \tiny\begin{ytableau}1&4\\2\\3\end{ytableau}。
\begin{itemize}
    \item \textbf{(12)}: $\text{diag}(1, -1, -1)$。
    \item \textbf{(23)}: 与 [3,1] 情形类似，$\rho=1/2$。涉及 $T_1, T_2$。
    \item \textbf{(34)}: 涉及 $T_2, T_3$。$3(1,2) \to 4(3,1)$。
    \item Path: Left(+1), Down(+1), Down(+1). Sum=3. $\rho=1/3$。
\end{itemize}
\noindent\textcolor{red}{\rule{\linewidth}{1pt}}
\section*{正多边形对称群 ($D_n$)}
\begin{itemize}
  \item \textbf{定义}：保持正$n$边形形状不变的对称变换群。
  \item \textbf{构成 ($2n$个元素)}：
  \begin{itemize}
      \item $n$个转动：绕$z$轴转角$2\pi/n$的整数倍。记$T=R_z(\frac{2\pi}{n})$，构成循环子群$Z_n=\{e, T, \dots, T^{n-1}\}$。
      \item $n$个翻转（二次固有转动）：绕平面内$n$条轴$S_j$转动$\pi$。
  \end{itemize}
  \item \textbf{二次轴$S_j$的几何分布}：
  \begin{itemize}
      \item $n$为偶数：$n/2$条连对边中点，$n/2$条连相对顶点。
      \item $n$为奇数：$n$条轴均为顶点与对边中点连线。
  \end{itemize}
  \item \textbf{乘法规则}（记轴$S_1,\dots,S_n$按逆时针排列）：
  \begin{itemize}
      \item 基本关系：$T^n=e,\quad S_j^2=e$。
      \item $TS_j = S_{j+1}$ (下标模$n$)，$S_jT = S_{j-1}$。
      \item 通式：$T^k S_j = S_{j+k},\quad S_j T^k = S_{j-k}$。
      \item 翻转之积：$S_{j+1} S_j = T,\quad S_j S_{j-k} = T^k$。
      \item 生成元：通常选$T$和$S_1$，则$S_k = T^{k-1}S_1$。
  \end{itemize}
  \item \textbf{共轭性质}：
  \begin{itemize}
      \item $T^k \sim T^{-k}$ 总是互相共轭。
      \item $S_i \sim S_{i+2}$ 总是互相共轭。
  \end{itemize}
  \item \textbf{共轭类分类 ($n=2k+1$ 奇数)}：共 $k+2$ 类。
  \begin{itemize}
      \item $\{e\}$：1类。
      \item $\{S_1, \dots, S_n\}$：所有二次轴共轭，1类。
      \item $\{T^r, T^{-r}\}$ ($1\le r \le k$)：$k$类。
  \end{itemize}
  \item \textbf{共轭类分类 ($n=2k$ 偶数)}：共 $k+3$ 类。
  \begin{itemize}
      \item $\{e\}$：1类。
      \item $\{T^k\}$：即转动$\pi$，与$T^{-k}$重合，自成1类。
      \item $\{S_{odd}\}, \{S_{even}\}$：奇数轴与偶数轴互不共轭，2类。
      \item $\{T^r, T^{-r}\}$ ($1\le r < k$)：$k-1$类。
  \end{itemize}
  \item \textbf{特例 $D_2$ 群}：
  \begin{itemize}
      \item 定义：$n=2$，元素$\{e, T, S_1, S_2\}$。
      \item 性质：Abel群，$T^2=S_1^2=S_2^2=e$。
      \item 关系：$TS_1=S_2, S_1S_2=T$ 等（地位全对称）。
      \item 二维表示：$e=(\begin{smallmatrix}1&0\\0&1\end{smallmatrix}), T=(\begin{smallmatrix}-1&0\\0&-1\end{smallmatrix}), S_1=(\begin{smallmatrix}-1&0\\0&1\end{smallmatrix}), S_2=(\begin{smallmatrix}1&0\\0&-1\end{smallmatrix})$。
  \end{itemize}
\end{itemize}
\noindent\textcolor{red}{\rule{\linewidth}{1pt}}
\section*{正多面体 (Regular Polyhedra)}
\begin{itemize}
  \item \textbf{定义参数}：
  $N$：面数 (Faces)；
  $V$：顶点数 (Vertices)；
  $L$：棱/边数 (Lines/Edges)；
  $n$：每个面是正$n$边形；
  $m$：每个顶点连接$m$条棱（也是$m$个面）。
  \item \textbf{拓扑关系}：
  \begin{enumerate}
      \item \textbf{计数关系}：$nN = 2L = mV$
      \item \textbf{欧拉公式 (Euler Characteristic)}：$V - L + N = 2$
  \end{enumerate}
\end{itemize}

\section*{存在性约束}
\begin{itemize}
  \item \textbf{几何角度}：正$n$边形内角 $\alpha = \frac{(n-2)\pi}{n}$。
  \item \textbf{立体角条件}：为构成闭合立体，顶点处棱数 $m \ge 3$，且内角和 $m\alpha < 2\pi$。
  \item \textbf{不等式判据}：
  \[
  3 \le m < \frac{2n}{n-2}
  \]
  \item \textbf{推论}：$n \ge 6$ 时无解（平面或双曲），仅 $n=3,4,5$ 有解。
\end{itemize}

\section*{正多面体分类表}
自然界仅存在以下5种正多面体（柏拉图立体）：
\begin{center}
\setlength{\tabcolsep}{1pt}
\begin{tabular}{|c|cc|ccc|}
\hline
\textbf{名称} & \textbf{面$n$} & \textbf{顶$m$} & \textbf{面$N$} & \textbf{顶$V$} & \textbf{棱$L$} \\
\hline
正4面体 & 3 & 3 & 4 & 4 & 6 \\
正6面体 (立方体) & 4 & 3 & 6 & 8 & 12 \\
正8面体 & 3 & 4 & 8 & 6 & 12 \\
正12面体 & 5 & 3 & 12 & 20 & 30 \\
正20面体 & 3 & 5 & 20 & 12 & 30 \\
\hline
\end{tabular}
\end{center}

\section*{正多面体的对偶 (Duality)}
\begin{itemize}
  \item \textbf{定义}：互换面与顶点（$N \leftrightarrow V$），棱数$L$保持不变。
  \item \textbf{对偶对}：
  \begin{itemize}
      \item 正六面体 $\longleftrightarrow$ 正八面体
      \item 正十二面体 $\longleftrightarrow$ 正二十面体
      \item 正四面体 $\longleftrightarrow$ 正四面体（自对偶）
  \end{itemize}
\end{itemize}
\end{multicols*}
\end{document}

